\unit{강연 182 — 형식기하학과 대수기하학}
\sid{182}
\src{193}
\seminar{(1959년 5월)}
\paperhead{형식기하학과 대수기하학}%
  {알렉상드르 그로텐디크 지음}

\fsection{1}{스킴}

체 $k$ 위에 정의된 아핀 대수공간은 그 아핀 대수 $A$($k$ 위에 정의된
정칙함수들의 환)에 의해 본질적으로 결정된다는 것은 잘 알려져 있다.
대수공간의 사상 $X\longrightarrow Y$는 $k$-대수 준동형
$A(Y)\longrightarrow A(X)$와 일대일로 대응한다. 한 대수공간에 대응하는
아핀 대수는 유한 생성 $k$-대수이며, ``고전적'' 관점에서는 멱영원을 갖지
않는다. 역으로 이러한 유형의 모든 대수는 $k$ 위에 정의된 한 대수공간의
아핀 대수로 얻어진다. 따라서 아핀 대수공간에 관한 상황을 가환대수의
말로 해석하게 하는 잘 알려진 사전이 있다. 그렇게 하면 더 일반적인
명제들을 얻게 된다는 사실은 오래전부터 알려져 있었다. 실제로 문제의
환들이 방금 살펴본 유형이라고 가정할 필요가 대개 없어졌으며, 흔히 뇌터
가정이면 충분했기 때문이다. 특히 주어진 바탕체가 있든 없든, 이 환들이
멱영원을 포함하는 경우를 배제할 이유는 없었다. 지금까지 기하학자들은
이러한 암시를 받아들이기를 거부하고 멱영원을 갖지 않는 아핀 대수, 즉
구조층에 멱영원이 없는 대수공간(더구나 대개는 ``절대 기약''인
대수공간)만을 고집해 왔다. 강연자는 이러한 사고방식이 대수기하학에서
참으로 자연스러운 방법들의 발전에 심각한 장애가 되었다고 생각한다.

$A$를 가환환이라 하자. $A$의 소아이디얼들의 집합
$X=\operatorname{Spec}(A)$에는 자연스러운 위상, 곧 ``자리스키 위상''
또는 스펙트럼 위상이 주어진다는 것은 잘 알려져 있다. 한편 $X$ 위에는
가환환의 층 $\Osh_X$가 있으며, $\mathfrak p\in X$에서 그 줄기는
국소화환 $A_{\mathfrak p}$이고 그 단면환은 $A$와 동일시된다. 이로써
$X$는 \emph{환 달린 공간}이 되며, 이를 $A$의 \emph{소 스펙트럼}이라
한다. 환 준동형 $f:A\longrightarrow B$는 환 달린 공간의 사상
$f':\operatorname{Spec}(B)\longrightarrow\operatorname{Spec}(A)$를
정의하며, 그 바탕 집합 사상은 다름 아닌
$\mathfrak p\longmapsto f^{-1}(\mathfrak p)$이다. 이런 방식으로 얻어지는

\src{194}
$\operatorname{Spec}(B)$에서 $\operatorname{Spec}(A)$로 가는 환 달린
공간의 사상들은 준동형 $\Osh_x\longrightarrow\Osh_y$ ($x=f'(y)$)가
국소 준동형인 것들, 그리고 오직 그런 것들뿐이다(즉 극대 아이디얼의
역상은 극대 아이디얼이다).

$\operatorname{Spec}(A)$와 동형인 환 달린 공간을 \emph{아핀 스킴}이라 하고,
국소적으로 아핀인 환 달린 공간, 즉 모든 점이 유도된 구조에 관하여 아핀
스킴인 열린 근방을 갖는 환 달린 공간을 \emph{준스킴}이라 한다. 준스킴의
\emph{사상}은 자명한 방식으로 정의하며, 국소적으로 환 준동형에 대응한다.

준스킴 $S$를 고정하고 준스킴의 사상 $X\longrightarrow S$를 생각하면,
$S$는 기초체나 기초환(더 정확히는 섬유화의 밑공간)의 역할을 한다. 이때
$X$를 $S$-준스킴이라 한다. $S=\operatorname{Spec}(A)$이면 이는
$\Osh_X$가 $A$-대수의 층이라는 뜻이기도 하다. 따라서 모든 준스킴은
유일한 방식으로 $\mathbf Z$-준스킴으로 볼 수 있다. 물론 $S$-준스킴들은
범주를 이루며, 더 나아가 이 범주에서 두 대상 $X,Y$의 곱은 항상 존재하고
$X\times_S Y$로 나타낸다. 이 곱 개념으로 사상
$S'\longrightarrow S$에 대응하는 $S$-준스킴의 밑준스킴 변경을 정의할
수 있다. 실제로 $X\times_S S'$은 $S'$-준스킴으로 볼 수 있다.

$X\times_S X$의 대각선이 닫혀 있으면 $X$가 $S$ \emph{위에서 분리}되어
있다고 한다. $\mathbf Z$ 위에서 분리된 준스킴을 \emph{스킴}이라 한다.
그러면 이는 어떤 밑 대상 위에서도 분리되어 있다. 단순히 하기 위해 이제부터
스킴만을 다루고, 더 나아가 이들이 \emph{뇌터}라고 가정한다. 즉 이들은
뇌터 환의 스펙트럼인 아핀 열린집합들의 유한 합집합이다. $S$의 모든 아핀
열린집합 $U$에 대하여 $X$에서의 그 역상이 유한 개의 아핀 열린집합의
합집합이고, 이 열린집합들의 환이 $U$의 환 위의 유한 생성 대수이면 $X$가
$S$ \emph{위에서 유한형}이라고 한다. 이러한 $S$-스킴들이 본래 의미의
기하학적 연구에 알맞다. 특히 모든 $s\in S$에 대하여 $s$ 위에서 $X$의
섬유 $f^{-1}(s)$는 $S$에서 $s$의 국소환 $\Osh_s$의 잉여체
$\kappa(s)$ 위의 대수적 스킴이다. 그러므로 $X$는 어느 정도는 매개변수
$s$가 $S$를 움직이는 ``대수공간'' $f^{-1}(s)$들의 족으로 볼 수 있다.
(국소적 관점에서 $s$는 주어진 환의 소아이디얼들을 움직인다.) 물론
$\kappa(s)$들은 서로 다른 표수를 가질 수 있다. $S=\operatorname{Spec}(k)$이고
$k$가 체이면, 이제는 구조층이 멱영원을 가질 수 있다는 차이만 빼고 통상적인
``대수공간'' 개념을 본질적으로 되찾는다.

\src{195}
잘 알려진 개념들에서 착안하여 \emph{사영 사상}, 더 일반적으로는
\emph{고유 사상}의 개념을 정의한다. 이러한 사상은 유한형이다. 또한 닫힌
부분집합을 닫힌 부분집합으로 보내며, 임의의 밑준스킴 변경 뒤에도 이 성질을
유지한다.

$X$가 스킴(언제나 그렇듯 뇌터 스킴)이면, 층 $\Osh_X$는 [2]의 의미에서
\emph{연접인 환의 층}이다. 따라서 $X$ 위의 연접 가군층은 국소적으로 어떤
준동형 $\Osh_X^{m}\longrightarrow\Osh_X^{n}$의 여핵과 동형인 층들이기도
하다.

\fsection{2}{형식 스킴}

$X$를 스킴, $X'$을 $X$의 닫힌 부분집합이라 하자. 그러면
$X'=\operatorname{Supp}(\Osh_X/J)$인 $\Osh_X$의 연접 부분층 $J$가
존재한다(더구나 그러한 것 가운데 가장 큰 것도 존재한다). $X'$에
$\Osh_X/J$를 갖추면 $X_0$로 나타내는 스킴이 되며, 이러한 스킴을
``$X$의 닫힌 부분스킴''이라 한다. 또한 모든 $n$에 대하여 $X'$에
$\Osh_X/J^{n+1}$를 갖춘 것을 생각하고 $X_n$으로 나타낼 수 있다. 이는
바탕 집합은 여전히 $X'$이지만 구조층
$\Osh_{X_n}=\Osh_X/J^{n+1}$이 다른 $X$의 닫힌 부분준스킴이다. 분명히
$\Osh_{X_n}$들은 $X'$ 위의 환의 층들의 사영계를 이루며, 그 사영극한
$\overline{\Osh}_X$를 \emph{$X'$을 따른 $\Osh_X$의 형식적 완비화}라
한다. 이 환의 층을 갖춘 $X'$을 \emph{$X'$을 따른 $X$의 형식적
완비화}라 한다. 따라서 이는 환 달린 공간이지만 일반적으로 스킴은 아니다.
$X$ 위의 모든 연접층 $F$에 대하여 마찬가지로 $X'$을 따른 $F$의 완비화
\[
\overline F=\varprojlim_n F_n,
\qquad F_n=F\otimes_{\Osh_X}\Osh_X/J^{n+1}
\]
를 생각할 수 있다. 이는 $\overline X$ 위의 가군층이다. 그 단면들을
``$X'$을 따른 $F$의 형식적 단면''이라 하며,
$\varprojlim_n\Gamma(X',F_n)$의 원소들과 동일시한다. $F=\Osh_X$이면
자리스키의 의미에서 ``$X'$을 따른 $X$의 정칙함수''를 얻는다. 다만 이
용어는 고전적 용어와 충돌하므로 여기서는 따르지 않겠다.

\emph{형식 스킴}(뇌터라고 이해한다)이란 다음 조건을 만족하는 위상환의 층
$\Osh_{\mathfrak X}$가 주어진 위상공간 $\mathfrak X$를 말한다. 위상환의
층들의 동형
\[
\Osh_{\mathfrak X}=\varprojlim_n\Osh_n
\]
이 있고, 여기서 $\Osh_n$들은 $\mathfrak X$ 위의 환층들의 사영계를
이루어 각각 $\mathfrak X$를 스킴 $\mathfrak X_n$으로 만든다. 또한
$m\geq n$일 때 준동형 $\Osh_m\longrightarrow\Osh_n$은 전사이고 그 핵은
$J_m^{n+1}$이며, 여기서 $J_m$은

\src{196}
$\Osh_m\longrightarrow\Osh_0$의 핵이다. 그러면 $\Osh_{\mathfrak X}$가
연접이고 줄기들이 뇌터 국소환인 환의 층임을 보일 수 있다.

정의에 따라 앞에서와 같은 형식적 완비화 $\overline X$는 형식 스킴이고,
역으로 모든 형식 스킴은 \emph{국소적으로} 이 유형이다. 실제로 \emph{아핀}
형식 스킴(즉 $\mathfrak X_0$가 아핀이며, 따라서 모든 $\mathfrak X_n$이
아핀인 형식 스킴)을 주는 것은 $J$-아딕이고 분리·완비인 뇌터 위상환을
주는 것과 동치이다.

통상적인 스킴에 대한 사상, 유한형 사상, 고유 사상 등의 정의는 어려움
없이 형식 스킴으로 확장된다.

\fsection{3}{세 기본정리}

$f:X\longrightarrow Y$를 스킴들의 고유 사상(언제나 그렇듯 뇌터
스킴들의 사상), $Y'$을 $Y$의 닫힌 부분집합, $X'$을 $X$에서의 그
역상이라 하고, 이에 대응하는 형식적 완비화 $\overline Y$와
$\overline X$를 생각하자. 그러면 $f$는 형식 스킴들의 사상
$\overline f:\overline X\longrightarrow\overline Y$를 유도하며, 이
사상도 고유이다. $F$가 $X$ 위의 연접층이면 $\overline F$는
$\overline X$ 위의 연접층이다. 정리 1에서는 $X,Y,F$를 잊고 형식
스킴들의 고유 사상 $\overline f$와 $\overline X$ 위의 연접층
$\overline F$만을 본다. (다만 강연자가 완전한 증명을 써 둔 것은
$X,Y,f,F$에서 출발하는 경우뿐이다.)

\result{정리}{1 (유한성 정리)}
\begin{enumerate}
  \item[\textit{i.}] $R^{q}\overline f_{*}(\overline F)$는
    $\overline Y$ 위의 연접층이다.
  \item[\textit{ii.}] 자연스러운 준동형
    \[
      R^{q}\overline f_{*}(\overline F)
      \longrightarrow \varprojlim_n R^{q}f_{n*}(F_n)
    \]
    은 동형이다.
\end{enumerate}

이 명제에서는 $Y'$을 정의하는 $\Osh_Y$의 연접 부분층 $J$를 하나
선택했다고 가정한다. 그 역상으로 $X'$을 정의하는 $\Osh_X$의 연접
부분층을 얻고, 따라서 제2절에서처럼 $F_n,X_n,Y_n$과
$f_n:X_n\longrightarrow Y_n$을 정의한다. 임의의 두 형식 스킴 사이의
고유 사상에서 출발했다면 이 표기들을 명시하는 데 필요한 사소한 변경은
자명하다.

정리 1은 ``형식 코호몰로지''만을 다루었다. 다음 정리는

\src{197}
이를 ``대수적 코호몰로지''와 관련짓고, 대수적 코호몰로지와 해석적
코호몰로지의 비교에 관한 세르의 잘 알려진 정리 [4]와 유사하다.

\result{정리}{2 (제1비교정리)} $R^{q}f_{*}(F)$는 $Y$ 위의
연접층이며(이는 정리 1의 특수한 경우이다), 자연스러운 준동형
\[
  \overline{R^{q}f_{*}(F)}
  \longrightarrow \varprojlim_n R^{q}f_{n*}(F_n)
\]
은 동형이다.

\result{따름정리}{1} 다음 표준 동형들이 있다.
\[
  \overline{R^{q}f_{*}(F)}=R^{q}\overline f_{*}(\overline F).
\]

$q=0$일 때 이 따름정리는 자리스키의 ``정칙함수의 기본정리''를 일반화한
것이며, 우리는 여기서 자리스키의 ``연결성 정리''의 일반화를 이끌어 낼
것이다. 한편 $q=0$일 때 정리 1의 \textit{(ii)}는 자명하지만 정리 2나
그와 동치인 정식화인 따름정리 1은 전혀 그렇지 않음에 유의하자. 실제로
증명은 $q$에 대한 내림 귀납법으로 진행된다. $q$가 크면 양변이 0이므로
자명하고, 따라서 $q=0$인 경우는 귀납법의 마지막 단계, 말하자면 ``가장
어려운'' 경우로 나타난다.

\result{따름정리}{2} $Y=\operatorname{Spec}(A)$이고 $Y'$은 $A$의 아이디얼
$J$로 정의된다고 하자. 그러면 $X$ 위의 모든 연접층 $F$에 대하여
$H^{q}(X,F)$는 유한 생성 $A$-가군이며, 그 $J$-아딕 완비화는
$H^{q}(\overline X,\overline F)$이다.

끝으로 이 따름정리를
$H=\Homsh_{\Osh_X}(F,G)$에 적용하면 다음을 얻는다.

\result{따름정리}{3} $Y=\operatorname{Spec}(A)$이고 $Y'$은 $A$의 아이디얼
$J$로 정의된다고 하자. $F,G$가 $X$ 위의 두 연접층이면
$\operatorname{Hom}(F,G)$는 유한 생성 $A$-가군이고, 그 $J$-아딕
완비화는 $\operatorname{Hom}(\overline F,\overline G)$와 동일시된다.

물론 자연스러운 사상
$\operatorname{Hom}(F,G)\longrightarrow
\operatorname{Hom}(\overline F,\overline G)$은 준동형
$u:F\longrightarrow G$에 그 ``연속성에 의한'' 연장
$\overline u:\overline F\longrightarrow\overline G$를 대응시키는
사상이다(이 규약 아래에서 $\overline F$는 $F$에 관한 함자가 된다).

이제 $A$가 그 $J$-아딕 위상에 관하여 분리·완비라고 하자. 그러면 앞의
따름정리 2와 3에서 다음을 얻는다.
\[
  H^{q}(X,F)=H^{q}(\overline X,\overline F),
  \qquad
  \operatorname{Hom}(F,G)
  =\operatorname{Hom}(\overline F,\overline G).
\]

\src{198}
이 마지막 항등식은 $X$ 위의 연접층의 범주가 $\overline X$ 위의 연접층의
범주의 한 \emph{부분범주}(그 사상들은 유도된 사상들이다)와 동일시됨을
보여 준다. 사실 다음의 더 강한 명제가 성립한다.

\result{정리}{3} $\overline X$ 위의 가군층이 연접일 필요충분조건은
$X$ 위의 어떤 연접층 $F$에 대하여 $\overline F$ 꼴의 층과 동형인
것이다(정리 2의 따름정리 3에 따라 $F$는 정준 동형까지 정해진다).
[이제 $Y=\operatorname{Spec}(A)$이며, $A$는 $J$-아딕이고 분리·완비인
뇌터 위상환임을 상기하자.]

\result{따름정리}{1} $X$의 닫힌 부분스킴들은 $\overline X$의 닫힌 형식적
부분스킴들과 일대일로 대응한다.

실제로 이들은 각각 $\Osh_X$와 $\overline{\Osh}_X$의 연접 부분층들과
대응한다. 사상의 그래프를 닫힌 부분스킴으로 보면, 따름정리 1에서 다음을
얻는다.

\result{따름정리}{2} $X,Z$가 $A$ 위에서 고유인 두 스킴이라고 하자
($A$는 $J$-아딕이고 분리·완비인 뇌터 환이다). 그러면
$g\longmapsto\overline g$는 $X$에서 $Z$로 가는 $Y$-사상들과
$\overline X$에서 $\overline Z$로 가는 $\overline Y$-사상들 사이의
일대일 대응을 정의한다.

다시 말해 $A$ 위에서 고유인 대수적 스킴들은 $\overline Y$ 위에서 고유인
형식적 스킴들의 범주의 한 \emph{부분범주}(그 사상들은 유도된 사상들이다)로
나타난다. 그러나 $\overline Y$ 위에서 고유이면서 ``대수화 가능''하지
않은 형식적 스킴, 즉 $A$ 위에서 고유인 어떤 $X$에 대해서도 $\overline X$와
동형이 아닌 형식적 스킴이 존재한다는 점에 주의해야 한다(복소수체 위에
정의된 대수다양체에서 오지 않는 콤팩트 복소 해석다양체가 존재하는 것과
마찬가지이다). 이러한 형식적 스킴은 ``모듈라이 이론''에서 자연스럽게
나타난다. 다만 형식적 스킴이 대수화 가능한 다음의 흥미로운 특수한 경우를
기록해 두자.

\result{정리}{4} $A$를 잉여체 $k$를 갖는 완비 뇌터 국소환이라 하고,
$\mathfrak X$를 $A$ 위에서 고유인 형식적 스킴이라 하자(그
$\mathfrak m(A)$-아딕 위상을 갖춘다). 다음을 가정한다.
\begin{enumerate}
  \item[\textit{i.}] $\Osh_{\mathfrak X}$의 국소환들은 평탄
    $A$-가군이다. 동치로, $\Osh_{\mathfrak X}$와 $A$에 $A$의 극대
    아이디얼의 거듭제곱들이 정의하는 여과를 주면, 연관 등급 대상들 사이에
    다음 관계가 성립한다.
    \[
      \operatorname{gr}(\Osh_{\mathfrak X})
      \simeq
      \operatorname{gr}^{0}(\Osh_{\mathfrak X})
      \otimes_k\operatorname{gr}(A).
    \]
  \item[\textit{ii.}]
    $\mathfrak X_0=\mathfrak X\otimes_A k$를 $k$ 위의 대수적 스킴으로 보면
    \[
      H^2(\mathfrak X_0,\Osh_{\mathfrak X_0})=0
    \]
    이다.
  \item[\textit{iii.}] $\mathfrak X_0$는 사영적이다.
\end{enumerate}
이 조건들 아래에서 $\mathfrak X$는 대수화 가능하며, 더 정확히는 사영
$A$-스킴 $X$의 형식적 완비화 $\overline X$와 동형이다.

조건 \textit{(ii)}와 \textit{(iii)}는 특히 $\mathfrak X_0$가 $k$ 위의
비특이 곡선이면 성립한다. 따라서 정리 4는 특히 주어진 종수의 곡선에 관한
``모듈라이 이론''에 적용할 수 있다\ldots{} 정리 4를 어떻게 증명하는지
설명하자. 조건 \textit{(i)}와 \textit{(ii)}가 다음 사실을 함의함을
보인다. $\mathfrak X_0$ 위의 모든 연접층 가운데 국소적으로 구조층과
동형인 것은 $\mathfrak X$ 위의 같은 종류의 층을 환원하여 얻어진다
(아래의 명제 3 참조).
이제 $\mathfrak X_0$ 위의 ``풍부한'' 층에서 출발하여(\textit{(iii)}에
따라 그러한 층이 존재한다) 이를 $\mathfrak X$ 위의 가역층으로 올린다.
그런 다음 정리 1을 사용하여, 이 가역층의 어떤 양의 텐서 거듭제곱이
$\mathfrak X$에서 스킴 $\mathbf P_A^r$의 형식적 완비화로 가는 몰입
사상을 정의함을 증명한다($\mathbf P_A^r$은 $A$ 위의 차원 $r$인
``표준 사영공간''이다).

정리 1부터 3까지의 증명은 [1]을 참조하라.

\src{199}
\fsection{4}{연결성 정리와 자리스키의 ``주정리''에의 응용}

$f:X\longrightarrow Y$를 스킴들의 고유 사상이라 하자. 유한성 정리에 따라
$f_*(\Osh_X)=\mathcal A$는 $Y$ 위의 연접층이다. 더구나 이는 가환 대수의
층이므로 $Y$ 위에서 유한인 $Y$-스킴 $g:Y'\longrightarrow Y$에 대응한다.
이 스킴은 $Y$ 위에서 아핀이라는 조건, 즉 아핀 열린집합의 역상이 아핀이라는
조건과 $g_*(\Osh_{Y'})=\mathcal A$라는 조건으로 정의된다. 이제 $f$가
표준적으로 $f=gf'$으로 인수분해됨은 바로 알 수 있다. 여기서
$f':X\longrightarrow Y'$은 $X$에서 $Y'$으로 가는 사상이며
$f'_*(\Osh_X)=\Osh_{Y'}$을 만족한다. 이 $f$의 인수분해를 $f$의
\emph{스타인 인수분해}라 한다. 제1비교정리의 따름정리 1을 $f'$과
$Y'$에서 점 $y'$ 하나로 이루어진 부분에 적용하면
$f'^{-1}(y')=X'$이 연결임을 얻는다. 실제로 그렇지 않다면 $X'$을 따른
$X$의 형식적 단면들이 국소환을 이루지 않을 터인데,
$f'_*(\Osh_X)_{y'}=\Osh_{Y',y'}$의 완비화는 국소환이기 때문이다! 따라서
다음을 증명하였다.

\result{정리}{5 (자리스키의 ``연결성 정리'')} $f:X\longrightarrow Y$를
고유 사상이라 하자. 그러면 $f$는 동형을 제외하고 유일하게 $f=gf'$으로
인수분해된다. 여기서 $g:Y'\longrightarrow Y$는 유한 사상이고
$f':X\longrightarrow Y'$은 $f'_*(\Osh_X)=\Osh_{Y'}$을 만족한다(따라서
$g_*(\Osh_{Y'})=f_*(\Osh_X)$이다). $f'$의 섬유들은 연결이다. 즉 $f$의
섬유 $f^{-1}(y)$의 연결성분들의 집합은 $y$ 위에 놓인 $Y'$의 점들의
집합, 다시 말해 $f_*(\Osh_X)_y$의 극대 아이디얼들의 집합과 일대일로
대응한다.

여기서 연결성 정리의 통상적인 변형들을 곧바로 얻는다. 다음의 결과만
기록하자.

\src{200}
\result{따름정리}{1} $X$의 점 $x$가 그 섬유 $f^{-1}(y)$에서 고립되어
있을 필요충분조건은 섬유 $f'^{-1}(y')$($y'=f'(x)$)이 $x$ 하나로
이루어져 있는 것이다. 동치로, $f'$이 $x$의 한 근방에서 $y'$의 한
근방으로 가는 동형을 유도하는 것이다. 이러한 점들의 집합은 열린집합
$U$이고, $f'$은 $U$에서 $Y'$의 한 열린집합으로 가는 동형을 유도한다.

$f'$이 $x$에서 국소 동형사상임을 보이기 위하여, $f'$이 동형
$\Osh_{y'}\longrightarrow\Osh_x$를 유도함을 주목하자. 이는
$f'_*(\Osh_X)=\Osh_{Y'}$와 다음 사실에서 알 수 있다. 즉 $f'$은 닫힌
사상이고 $y'$ 위의 그 섬유는 $x$ 하나로 이루어져 있으므로, $V$가
$y'$의 근방들의 기본계를 움직일 때 $f'^{-1}(V)$는 $x$의 근방들의
기본계를 움직인다. 이로부터 ``기하학적'' 경우에는 슈발레에게서 비롯된
다음 결과를 곧바로 얻는다.

\result{따름정리}{2} $f$가 유한 사상일 필요충분조건은 $f$가 고유이고
그 섬유들이 유한한 것이다.

실제로 이 경우에는 앞의 결과에 따라 $f'$이 동형이다.

이제 $f:X\longrightarrow Y$를 반드시 고유일 필요는 없는 사상이라 하자.
다만 $X$가 고유 $Y$-스킴
$\widetilde f:\widetilde X\longrightarrow Y$의 열린 부분집합으로 들어
있다고 가정하자(특히 $f$가 준사영이면 그러하다). 따름정리 1을 적용하면,
$\widetilde f'$이 자기 섬유에서 고립된 $X$의 점들의 집합 $U$에서
$Y'$의 한 열린집합으로 가는 동형을 유도함을 알 수 있다(따라서 $U$는
실제로 열린집합이다). 이로부터 자리스키의 ``주정리''의 다음 전역 형태를
얻는다.

\result{정리}{6 (자리스키의 ``주정리'')} $f:X\longrightarrow Y$를
유한형 사상이라 하자. 그러면 자신이 속한 섬유에서 고립된 $X$의 점들의 집합
$U$는 열린집합이다. 또한 $f$가 분리 사상이면,\footnote{1962년의 정오표는
여기서 인쇄본의 ``$f$는 준사영이다''라는 가정을 더 약한 ``$f$는 분리
사상이다''라는 가정으로 바꾼다. 이는 SGA VIII, 6.2, 즉 모든 준유한 분리
사상은 준사영이라는 사실에 따른 것이다.} $U$는 $Y$ 위에서 유한인 어떤
스킴 $Y'$의 열린 부분과 $Y$-동형이다.

유한형 사상은 국소적으로 아핀이며, 따라서 \textit{a fortiori}
국소적으로 준사영이므로, 정리 6에서 주정리의 통상적인 국소 형태들이
곧바로 따른다.
\fsection{5}{고유 평탄 사상의 코호몰로지적 연구에의 응용}

$f:X\longrightarrow Y$를 고유 사상이라 하고, $F$를 $X$ 위의 연접층이라
하자. $F$는 $Y$ 위에서 평탄하다고 가정한다. 즉 $F_x$들은 환
$\Osh_y$ 위의 평탄 가군이다($y=f(x)$). 이는 또한 다음을 뜻한다. 모든
$y\in Y$에 대하여 섬유 $f^{-1}(y)$을 따라 $F$에
$\mathfrak m_y^nF$들로 이루어진 여과를 주면($\mathfrak m_y$는
$\Osh_y$의 극대 아이디얼이다), 그 연관 등급 대상은
$(F/\mathfrak m_yF)\otimes_{\kappa(y)}\operatorname{gr}(\Osh_y)$와
동형이다. 다시 말해 모든 정수 $n$에 대하여
\[
  \mathfrak m_y^nF/\mathfrak m_y^{n+1}F
  =F_y\otimes_{\kappa(y)}
  (\mathfrak m_y^n/\mathfrak m_y^{n+1})
\]
이다. 여기서 $X_y$는 섬유 $f^{-1}(y)$을 $y$의 잉여체 $\kappa(y)$
위에서 고유인 스킴으로 본 것이고, $F_y$는 $F$가 $X_y$ 위에 유도하는
층 $F/\mathfrak m_yF$이다. 이 동형과 정리 2를 고려하면, $F_y$를
계수로 하는 $X_y$의 코호몰로지를 앎으로써 $y$의 근방에서
$R^qf_*(F)$의 상계를 얻고 때로는 이를 계산할 수 있다. 여기서 정리 2는
다음 꼴이 된다.
\[
  \widehat{R^qf_*(F)_y}
  =\varprojlim_n H^q(X_y,F/\mathfrak m_y^nF).
\]
여기서는 그 따름정리 하나만 기록하겠다.

\src{201}
\result{명제}{1} $f:X\longrightarrow Y$를 고유 사상, $F$를 $X$ 위의
연접층이라 하고, $F$가 $Y$ 위에서 평탄하다고 하자. $y\in Y$이고 $q$가
정수이며 $H^q(X_y,F_y)=0$이라고 가정하자. 그러면 $R^qf_*(F)$는 $y$의
한 근방에서 영이고, 모든 $n$에 대하여 자연스러운 준동형
\[
  R^{q-1}f_*(F)_y\longrightarrow
  H^{q-1}(X_y,F/\mathfrak m_y^nF)
\]
은 전사이다.

특히 $f$가 평탄 사상이면(즉 $\Osh_X$가 $Y$ 위에서 평탄하면), $X$ 위의
모든 국소 자유 연접층 $F$는 $Y$ 위에서 평탄하다. 이러한 두 층 $F,G$에
대하여 명제 1을 $\Homsh_{\Osh_X}(F,G)$와 $q=1$에 적용하면 다음을 얻는다.

\result{정리}{7} $f$를 고유 평탄 사상, $F,G$를 $X$ 위의 국소 자유인 두
연접층, $y\in Y$라 하고
\[
  H^1\bigl(X_y,\Homsh_{\Osh_{X_y}}(F_y,G_y)\bigr)=0
\]
이라고 가정하자. 그러면 모든 준동형 $u_0:F_y\longrightarrow G_y$는
준동형
\[
  u:F|V\longrightarrow G|V,
  \qquad V=f^{-1}(U),
\]
에서 유도된다. 여기서 $U$는 $y$의 한 근방이다.

\result{따름정리}{} $u_0$가 동형사상(각각 단사사상, 에피사상)이면,
$U$를 충분히 작게 잡았을 때 $u$도 마찬가지이다.

특히 다음이 성립한다.

\src{202}
\result{따름정리}{2} $E_0$를 $X_y$ 위의 국소 자유 연접층이라 하고
\[
  H^1\bigl(X_y;\Homsh_{\Osh_{X_y}}(E_0,E_0)\bigr)=0
\]
이라고 하자. 그러면 $X_y$에 제한했을 때 $E_0$와 동형인 두 국소
자유층은 $X_y$의 한 근방에서 서로 동형이다.

따라서 다음이 성립한다.

\result{따름정리}{3} $H^1(X_y,\Osh_{X_y})=0$이라고 가정하자. 그러면
$X$ 위의 두 가역층(즉 $\Osh_X$와 국소적으로 동형인 층)은 $X_y$에 대한
그 제한들이 서로 동형이면 서로 동형이다.

이로부터 다음을 얻는다.

\result{명제}{2} $Y$를 연결 스킴, $E$를 $Y$ 위의 계수가 적어도 $2$인
국소 자유 연접층이라 하자. $E$로 정의되는 사영 다발
$X=\mathbf P(E)$를 그 잘 알려진 가역층 $\Osh_X(1)$과 함께 생각하자.
그러면 $X$ 위의 모든 가역층 $L$은
\[
  f^*(L')\otimes\Osh_X(n)
\]
꼴의 층과 동형이다. 여기서 $L'$은 $Y$ 위의 가역층이고 $n$은 정수이다.
이 $n$은 유일하게 정해지고, $L'$은 동형을 제외하고 정해진다.

앞의 따름정리 3에 따라 $L$은 각 섬유의 한 근방에서 어떤
$\Osh_X(n)$과 동형이다. 나머지는 거의 형식적으로 따른다.

명제 2를 사용하면 $X=\mathbf P(E)$에서 다른 사영 다발로 가는
$Y$-사상들을 결정할 수 있다. 특히 다음을 얻는다.

\result{따름정리}{} $u$를 $X=\mathbf P(E)$의 자동사상이라 하자. 그러면
$Y$ 위의 가역층 $L'$과 $E$에서 $E\otimes L'$으로 가는 동형사상 $v$가
존재하여, $u$는 이에 대응하는 동형사상
\[
  \mathbf P(E)\xrightarrow{\sim}\mathbf P(E\otimes L')=\mathbf P(E)
\]
이 된다. 쌍 $(v,L')$은 동형을 제외하고 정해진다.

$\Gamma$를 $E\otimes L'$이 $E$와 동형이 되는 $Y$ 위의 가역층 $L'$의
동형류들로 이루어진 집합이라 하자. 그 원소들은 모두 유한 위수를 갖는다.
실제로 $n$이 $E$의 계수이면, $n$차 외승을 취하여
$L'^{\otimes n}\simeq\Osh_Y$이어야 함을 알 수 있다. 따라서 이
따름정리는 다음 군의 완전열이 존재한다는 말로 표현할 수 있다.
\[
  e\longrightarrow
  \operatorname{Aut}(E)/\Gamma(Y,\Osh_Y^*)
  \longrightarrow \operatorname{Aut}_Y(X)
  \longrightarrow \Gamma\longrightarrow e
\]
(사실 이는 다음 군의 층들의 완전열에서 유도되는 코호몰로지 완전열로부터도
얻어진다.
\[
  e\longrightarrow \Osh_Y^*
  \longrightarrow \underline{\operatorname{Aut}}(E)
  \longrightarrow \underline{\operatorname{Aut}}_Y(X)
  \longrightarrow e,
\]
여기서 $\Osh_Y^*$는 $\Osh_Y$의 ``가역원''들의 층이며
$\underline{\operatorname{Aut}}(E)$의 중심과 동일시된다).

\src{203}
\fsection{6}{$J$-아딕 완비환 위의 층과 스킴에 대한 존재 및 유일성 정리의 응용}

정리 7에서는 정리 1과 2를 이용하여 국소 자유 연접층에 관한 유일성 결과를
주었다. 이제 정리 3을 이용하여 층, 스킴의 사상 또는 스킴 자체에 관한 존재
정리들을 도출한다. 이하에서 $A$는 분리·완비인 뇌터 국소환을 나타낸다.
일반적인 방법은 언제나 형식적 구성을 수행하는 데 있다. 이는 본질적으로
아르틴 환 위에서 대수기하학을 수행하고, 세 기본정리를 이용하여 그로부터
``대수적'' 성질의 결론을 끌어내는 것이다.

\result{명제}{3} $\mathfrak X$를 $A$ 위에서 고유이고 평탄인 형식적
스킴이라 하고, $F_0$를 $X_0$ 위의 국소 자유층이라 하자. 다음을 가정한다.
\[
  H^2\bigl(X_0,\Homsh_{\Osh_{X_0}}(F_0,F_0)\bigr)=0.
\]
그러면 $X_0$ 위에 $F_0$와 동형인 층을 유도하는 $\mathfrak X$ 위의 국소
자유층 $F$가 존재한다. (더구나
$H^1(X_0,\Homsh_{\Osh_{X_0}}(F_0,F_0))=0$이면 이 $F$는 동형까지
유일하다.)

$X_n$ 위의 국소 자유층 $F_n$들을 각각 앞의 것을 유도하도록 차례로
구성한다. $F_n$을 구성하는 데 따르는 장애는
\[
  H^2\bigl(X_0,\Homsh_{\Osh_{X_0}}(F_0,F_0)\bigr)
  \otimes_{A/J}(J^n/J^{n+1})
\]
안에 놓이므로, 가설에 따라 사라진다. 이제 정리 3을 이용하면 다음을 얻는다.

\result{따름정리}{1} $X$를 $A$ 위에서 고유이고 평탄인 스킴이라 하고,
$F_0$를 위와 같이 잡자. 그러면 $X_0$ 위에 $F_0$와 동형인 층을 유도하는
$X$ 위의 국소 자유층 $F$가 존재한다. 더구나
$H^1(X_0,\Homsh_{\Osh_{X_0}}(F_0,F_0))=0$이면 이 $F$는 동형까지
유일하다.

$X_0$를 체 $k$ 위의 유한형 스킴이라 하자. $X_0$는 $k$ 위에서
비특이(여기서는 절대 비특이라는 뜻이다)라고 가정하지만, 반드시 $k$ 위에서
고유일 필요는 없다. $A$를 잉여체가 $k$인 아르틴 국소환이라 하자. 우리는
$A$ 위에서 평탄이고
$X\otimes_A k=X_0$인 스킴 $X$를 찾는 데 관심이 있다(이것이 $X_0$의
``모듈라이 이론'', 또는 ``구조의 변형''의 출발점이다). 그러한 $X$를 주는
것과, 위상공간 $X_0$ 위의 층 $\Osh_X$에 다음 구조들을 주는 것은 같다.
\begin{enumerate}
  \item[\textit{i.}] 이는 $A$-대수의 층이다.
  \item[\textit{ii.}] $A$-대수 구조들과 양립하는 증대 준동형
    $\Osh_X\longrightarrow\Osh_{X_0}$을 갖는다. 이 데이터에는 다음 조건들이
    부과된다. 증대 준동형은 동형사상
    $\Osh_X\otimes_A k\xrightarrow{\sim}\Osh_{X_0}$을 유도한다. 또한
    $\Osh_X$는 $A$ 위에서 평탄하다. 즉 $A$의 극대 아이디얼
    $\mathfrak m$의 거듭제곱들로 여과한 $\Osh_X$의 연관 등급 대상은
    $\operatorname{gr}^0(\Osh_X)\otimes_k\operatorname{gr}(A)$와 동형이다.
    다시 말해 다음 동형사상들이 존재한다.
    \[
      \mathfrak m^n\Osh_X/\mathfrak m^{n+1}\Osh_X
      =\Osh_{X_0}\otimes_k(\mathfrak m^n/\mathfrak m^{n+1}).
    \]
\end{enumerate}
기본 사실은 다음과 같다.

\src{204}
\result{정리}{8} $X_0$를 체 $k$ 위의 유한형 비특이 스킴이라 하고,
$X_0$가 아핀이라고 가정하자. $A$를 잉여체가 $k$인 아르틴 국소환이라
하자. 그러면 $A$ 위에서 평탄이고 $X\otimes_A k=X_0$인 $A$-스킴 $X$가
존재하며, 그러한 두 스킴은 반드시 서로 동형이다.

여기서 말한 동형사상은 표준적이지 않음에 유의하자. 일반적으로 $X$에는
$X_0$ 위에서 항등을 유도하는 자명하지 않은 $A$-자기동형사상들이 있기
때문이다. 한편 주어진 조건들을 만족하는 $X$의 ``표준적'' 선택은 일반적으로
존재하지 않는다. 다만 $A$가 $k$-대수인 경우(동표수인 경우)에는
$X=X_0\otimes_k A$, 즉 $\Osh_X=\Osh_{X_0}\otimes_k A$로 잡을 수 있다
($X_0$가 아핀인지 여부와 무관하다). 이표수인 경우에는 $X_0$가 아핀이 아닐
때 $X_0$를 $A$ 위에서 정의된 $X$로 ``올릴'' 수 있는지 나는 일반적으로
알지 못한다. 그러나 $n>0$인 정수 $n$을 잡고
$A_{n-1}=A/\mathfrak m^n$으로 놓자. $X_0$를 평탄 $A_{n-1}$-스킴
$X_{n-1}$로 올렸다고 가정하고, $X_{n-1}$을 평탄 $A_n$-스킴 $X_n$으로
올리고자 한다. 정리 8에 따라 이것이 국소적으로 가능함은 이미 알고 있다.
한편 $U_n$이 $X_{n-1}$의 열린집합 $U_{n-1}$을 올린다면, $U_{n-1}$ 위에서
항등을 유도하는 $U_n$의 자기동형군의 층은
\[
  \Theta_{X_0/k}\otimes_k(\mathfrak m^n/\mathfrak m^{n+1})
\]
을 $U_{n-1}$에 제한한 층과 표준적으로 동형이며, 특히 가환임을 쉽게
확인할 수 있다. (여기서 $\Theta_{X_0/k}$는 $X_0$ 위의 $k$-미분들의
싹의 층을 나타낸다.) 따라서 $X_{n-1}$을 올리는 $X_n$을 구성하는 데
따르는 장애가 존재하며, 이는
\[
  H^2(X_0,\Theta_{X_0/k})\otimes_k
  \mathfrak m^n/\mathfrak m^{n+1}
\]
안에 놓인다. 따라서 다음을 얻는다.

\result{따름정리}{1} $X_0$를 $k$ 위의 유한형 비특이 스킴이라 하고,
$H^2(X_0,\Theta_{X_0/k})=0$이라고 가정하자. 그러면 잉여체가 $k$인 모든
아르틴 국소환 $A$에 대하여 $A$ 위에서 평탄이고
$X\otimes_A k=X_0$인 $A$-스킴 $X$가 존재한다.

더구나 $X_0$를 올리는, $A$ 위에서 평탄인 $X$를 하나 찾았다면, 정리 8에
따라 $X_0$를 올리는 평탄 $A$-스킴들의 동형류의 집합은
\[
  H^1\bigl(X_0,\underline{\operatorname{Aut}}(X)\bigr)
\]
과 동일시된다. 여기서 물론 $\underline{\operatorname{Aut}}(X)$는 증대와
양립하는 $A$-대수의 층 $\Osh_X$의 자기동형사상들의 싹의 층을 나타낸다.
$\Osh_X$의 여과는 $\underline{\operatorname{Aut}}(X)$의 여과를 정의하며,
이 층을 그 여과의 $n$번째 부분군으로 나눈 몫은
$\underline{\operatorname{Aut}}(X_n)$이다. 이 여과의 연관 등급 대상은
가환이고 $\Theta_{X_0/k}\otimes_k\operatorname{gr}(A)$와 동일시된다.
특히 $\mathfrak m^{n+1}$이 처음으로 0이 되는 $\mathfrak m$의 거듭제곱이면,
$F^n(\underline{\operatorname{Aut}}(X))$(여과의 마지막 단계)는
$\underline{\operatorname{Aut}}(X)$의 중심에 들어가며
$\Theta_{X_0/k}\otimes_k\mathfrak m^n$과 동형이다. 이것은 또한
$X_{n-1}=X\otimes_A A/\mathfrak m^n$ 위에서 항등을 유도하는 $X$의
자기동형사상들의 싹의 층이다. 이 결과들을 이용하면 다음 명제들을 곧바로
얻는다.

\src{205}
\result{따름정리}{2} $X_0$를 $k$ 위의 유한형 비특이 스킴이라 하고,
$A$를 잉여체 $k$와 극대 아이디얼
$\mathfrak m$을 갖는 아르틴 국소환이라 하자. $\mathfrak m^{n+1}=0$이라 가정하고,
$A_{n-1}=A/\mathfrak m^n$이라 하며, 끝으로 $X_{n-1}$을 평탄
$A_{n-1}$-스킴으로서 $X_{n-1}\otimes_A k=X_0$인 것이라 하자. 그러면
$X_{n-1}$ 위에서 항등사상을 유도하는 동형을 제외한, 평탄
$A$-스킴 $X_n$ 중 $X_n\otimes_A A_{n-1}=X_{n-1}$인 것들의
동형류 집합은 공집합이거나
다음 군에 대한 주동차공간이다.
\[
  H^1(X_0,\Theta_{X_0/k})\otimes_k\mathfrak m^n.
\]

(일반적으로 이 마지막 공간에는 특별히 정해진 원점이 없음을 유의하라.
$X_{n-1}$을 $X_n$으로 올리는 특별히 정해진 방법이 없기
때문이다.)

\result{따름정리}{3} $X_0$를 $k$ 위의 유한형 비특이 스킴이라 하고,
$H^1(X_0,\Theta_{X_0/k})=0$이라 가정하자. 그러면 잉여체 $k$를 갖는 모든
아르틴 국소환 $A$에 대하여, $X\otimes_A k=X_0$인 평탄
$A$-스킴 $X$는 동형을 제외하면 많아야 하나 존재한다.

따름정리 1과 3에서, $A$를 잉여체 $k$를 갖는 완비 뇌터
국소환으로만 가정하고 $X$를
$A$ 위의 형식적 스킴으로 취했을 때의, 겉보기에는 더 일반적인 명제들이 즉시 따른다.

\result{정리}{9} $k$를 체, $X_0$를 $k$ 위의 유한형 비특이 스킴이라
하자. 잉여체가 $k$인 모든 완비 뇌터 국소환 $A$에 대하여,
$F(A)$를 $X_0$ 위에서 항등사상을 유도하는 동형을 제외한, $A$ 위의
유한형 평탄 형식적 스킴 $\mathfrak X$ 중
$\mathfrak X\otimes_A k=X_0$인 것들의 동형류 집합이라 하자. 이 표기 아래에서 다음이 성립한다.
\begin{enumerate}
  \item[\textit{i.}] $H^1(X_0,\Theta_{X_0/k})=0$이면 모든 $A$에 대하여
    $F(A)$는 많아야 한 원소를 갖는다.
  \item[\textit{ii.}] $H^2(X_0,\Theta_{X_0/k})=0$이면 모든 $A$에 대하여
    $F(A)$는 적어도 한 원소를 갖는다.
\end{enumerate}

\src{206}
\result{따름정리}{1} $X_0$가 $k$ 위에서 고유라고 가정하자. 조건
\textit{(i)} 아래에서 모든 $A$에 대하여, $X_0$ 위에서 항등사상을
유도하는 동형을 제외하면, $X\otimes_A k=X_0$인
$A$ 위의 고유 평탄 스킴 $X$는 많아야 하나 존재한다.

여기서는 정리 3의 따름정리 2를 사용한다. 예를 들면 다음과 같다.

\result{따름정리}{2} $X$가 $A$ 위의 고유 평탄 스킴이고,
$X\otimes_A k$가 $k$ 위의 차원 $r$인 표준 사영공간
$\mathbf P_k^r$와 동형이면, $X$는 $\mathbf P_A^r$와 동형이다.

(이 결과는 명제 3의 따름정리 1에서도 도출할 수 있다.)

\result{따름정리}{3} $X_0$를 $k$ 위의 비특이 사영 스킴이라 하고,
다음을 가정하자.
\[
  H^2(X_0,\Osh_{X_0})=H^2(X_0,\Theta_{X_0/k})=0.
\]
그러면 모든 $A$에 대하여, $X\otimes_A k=X_0$인 사영
평탄 $A$-스킴 $X$가 존재한다.

정리 9\textit{(ii)}와 정리 4를 결합한다. 특히 다음을 얻는다.

\result{따름정리}{4} $X_0$를 $k$ 위의 완비 비특이 대수곡선의
스킴이라 하자. 그러면 잉여체가 $k$인 모든 완비 뇌터 국소환 $A$에 대하여,
$X\otimes_A k=X_0$인 ``비특이 곡선 스킴'' $X$가 $A$ 위에 존재한다.

\result{비고}{}

1) 따름정리 3과 4는 특히 $k$의 표수가
$p\ne0$일 때 흥미롭다. 이때 $A$로 표수가 $0$이고 잉여체가 $k$인
이산 값매김환을 취한다. 예를 들면 ``가능한 가장 작은 $A$'',
즉 $p$가 극대 아이디얼을 생성하는 환을 취한다. (실제로
Cohen의 정리들에 따라, 그러한 환 $A$ 하나에 대하여 따름정리 3과 4를
아는 것으로 충분하다.) 이와 관련하여 전문가들의 말에 따르면,
체 $k$ 위에 정의된 스킴 가운데 그러한 환 $A$ 위에 정의된 평탄 스킴의
mod $p$ 환원이 아닌 것이 존재하는지는 알려져 있지 않다. 적어도 이 절의 결과들은
이 문제를 체계적으로 조사할 방법을 준다. 먼저
$H^2(X_0,\Theta_{X_0/k})$에 놓이는 첫 번째 장애가
반드시 영인지 살펴보아야 할 것이다.

\emph{1962년 보충.} --- J.-P. Serre는 표수가 $p>0$인 대수적으로
닫힌 체 $k$ 위에 3차원 비특이 사영 다양체를 구성했는데, 이 다양체는
잉여체가 $k$이고 분수체의 표수가 영인 국소 정역 위의 고유 스킴을
환원하여 얻어지는 것이 아니다. Mumford는 비특이 사영 곡면에 관한
이와 비슷한 예를 찾았다고 한다.

2) 정리 3과 이에 대응하는 기법은 밑환을 (생각을 구체화하기 위해 국소환으로
잡을 때) 완비라고 해야만 유효함을 유의하자.

국소환의 완비화에 대하여 알려진 결과에서 그 국소환 자체에 대한
대응하는 결과로 넘어가려면, 그 최종적인 명제가 아직 찾아지지 않은
네 번째 ``기본정리''가 필요할 것이다.

\src{207}
3) 이 절의 결과들(특히 앞의 따름정리 1과 2)과 다음 절의 결과들을,
복소구조의 변형에 관한 Kodaira--Spencer의 결과들과 비교해야 한다.
방금 언급한 추측적 정리를 사용하면, 따름정리 1의 조건 아래에서
$A$가 더 이상 완비라고 가정되지 않더라도, $A$를 포함하면서
$A$ 위에서 유한이고 비분기인 환 $A'$이 존재하여
$X\otimes_A A'$과 $X'\otimes_A A'$이 $A'$-동형임을 결론지을 수 있어야 한다
($X,X'$은 $X\otimes_A k=X'\otimes_A k=X_0$을 만족하는, 주어진 두
고유 평탄 $A$-스킴이다). 적어도
$X_0=\mathbf P_k^r$일 때에는 명제 2의 따름정리를 사용하여 이를 증명할 수 있다. 어떤
경우든 $H^1(X_0,\Theta_{X_0/k})=0$이면, $Y=\operatorname{Spec}(A)$의
모든 점 $y$ 위에서 $X,X'$의 섬유들은
적어도 잉여체 $\kappa(y)$의 대수적 폐포로 넘어간 뒤에는 동형임을 증명할 수 있다.
(더 이상 $A$가 반드시 국소환이라고 가정하지 않는, 겉보기에는
더 강한 국소 결과도 있다.) 사영공간의 ``구조의 변형''과 관련하여,
Kodaira--Spencer의 대응하는 문제에서 제안된 다음 질문도 언급하자. $X$를
분수체 $K$와 잉여체 $k$를 갖는 국소 정역 $A$ 위의
고유 평탄 스킴이라 하자. $X\otimes_A K$가
$\mathbf P_K^r$와 동형이라고 가정할 때, $X\otimes_A k=X_0$는
$\mathbf P_k^r$와 동형인가(적어도 $k$의 대수적 폐포 위에서)? 이 문제에서는
$A$를 완비 이산 값매김환이라고 가정해도 된다. $X_0$가
아벨 다양체일 때에도 비슷한 문제가 있다.

\emph{1962년 보충.} --- 강연자는 이때 위에서 추측한 결과들에 대하여
덜 낙관적이다. 그러나 사영공간의 구조의 변형에 관한 문제는 Hironaka가
긍정적으로 해결했고, 아벨 다양체에 관한 비슷한 문제는 Koizumi가 해결했다.
\fsection{7}{``모듈라이 이론''에의 응용}

강연자는 이 이론을 이제 막 직접 다루기 시작했으므로, 여기서는 몇 가지
지적에 한정하지 않을 수 없다. 단순화를 위해 체 $k$ 위에 놓는다. 즉 우리는
동일 표수에서 작업한다. 정리 8을 사용하면 본질적인 변경 없이 일반적인
경우도 다룰 수 있을 듯하다. 현재 우리는 아직 ``형식적'' 단계를 넘어서지
못했다. 그러나 강연자는 여기에서 출발하여 몇몇 경우에는 실제 모듈라이
스킴을 구성하고, 특히 모든 정수 $g$에 대하여 종수 $g$인 비특이 곡선들의
보편 모듈라이 스킴 역할을 하는 정수환 위의 스킴을 구성할 수 있으리라
기대한다.

\emph{1962년 보충.} --- Mumford는 종수 $g$인 곡선들의 모듈라이 스킴을
구성하였다(cf. SHMI). 또한 정리 10은 모듈라이 이론의 ``레벨 $n$의 야코비
스킴''이 비특이이고, 더 나아가 $\mathbf Z$ 위에서 매끄럽다는 것을 증명한다.

\src{208}
이제 정리 9의 상황과 표기로 돌아가되, $A$가 $k$ 위의 유한 차원 국소대수라고
가정하자. 단순화를 위해 $k$는 대수적으로 닫혔다고 가정한다. 그러면 $F(A)$를
$A$에 관하여 집합의 범주에 값을 갖는 공변 함자로 볼 수 있다. 실제로
$k$-대수의 준동형 $A\longrightarrow B$는 사상
$F(A)\longrightarrow F(B)$를 정의한다. 왜냐하면
$X\otimes_A k=X_0$을 만족하는 모든 평탄 $A$-스킴 $X$는 같은 성질들을 갖는
$B$-스킴 $X\otimes_A B$를 낳기 때문이다. 완비 뇌터 국소 $k$-대수
$\mathcal O$와 함자적 동형사상
\begin{equation}
  \tag{*}
  \operatorname{Hom}(\mathcal O,A)\xrightarrow{\sim}F(A)
\end{equation}
을 찾을 수 있다고 가정하자(여기서 왼쪽 항은 $k$-대수의 준동형들을 뜻한다).
그러한 $\mathcal O$는 정준 동형까지 정해짐을 쉽게 알 수 있다. 이때
$\mathcal O$의 형식적 스펙트럼 $\mathfrak Y$, 즉 한 점으로 이루어진
위상공간에 $\mathcal O$ 하나로 이루어진 위상환의 층을 준 것을 $X_0$의
\emph{형식적 모듈라이 스킴}이라 한다(이것이 반드시 존재하는 것은 아니다).
$\mathfrak r$을 $\mathcal O$의 극대 아이디얼이라 하고, 모든 $n$에 대하여
$\mathcal O_n=\mathcal O/\mathfrak r^{n+1}$이라 하자(따라서
$\mathcal O_0=k$이다). 그러면 표준 준동형
$\mathcal O\longrightarrow\mathcal O_n$은
$\operatorname{Hom}(\mathcal O,\mathcal O_n)$의 원소이므로
$F(\mathcal O_n)$의 원소, 즉 $\mathfrak r$로 나눈 환원이 $X_0$인 평탄
$\mathcal O_n$-스킴 $X_n$을 정의한다. 이 $X_n$들은 스칼라 확대(즉 여기서는
환원)에 의해 서로에게서 얻어진다. 따라서 이들은 형식적 모듈라이 스킴
$\mathfrak Y$ 위의 잘 정해진 형식적 스킴 $\mathfrak X$에서 온다.
$\mathfrak X$는 $\mathfrak Y$ 위에서 평탄이고 $\mathfrak X_0=X_0$이다.
그러면 동형사상 (*)은, 바로 알 수 있듯이, 모든 $k$-대수의 준동형
$\mathcal O\longrightarrow A$에 $A$-스킴
$\mathfrak X\otimes_{\mathcal O}A$의 동치류를 대응시켜 주어진다. 즉 모든
$k$-스킴의 사상 $Y'=\operatorname{Spec}(A)\longrightarrow\mathfrak Y$에
밑준스킴의 변경으로 얻은 $Y'$-스킴
$\mathfrak X\otimes_{\mathfrak Y}Y'$을 대응시킨다. 더 나아가 $A$가 단지
완비 뇌터 국소 $k$-대수라고만 가정해도(반드시 아르틴일 필요는 없다)
동형사상 (*)과 방금의 설명은 여전히 성립함을 알 수 있다. 물론 통상 그렇듯
$\mathcal O$는 \textit{a priori} 멱영원을 가질 수 있으며,
$\mathcal O$ 자체가 $k$와 같지 않은 아르틴 환인 경우도 존재할 법하다. 이는
$A$를 정칙환으로 한정하는 Kodaira--Spencer의 관점이 일반적인 경우에는
\textit{a priori} 얼마나 부적절한지를 보여 준다.

\src{209}
$k$ 위에서 고유라고 가정한 $X_0$에 형식적 모듈라이 스킴이 존재하기 위한
충분조건을 제시하는 일이 남았다. 일반적으로, 유한 차원 국소 $k$-대수에서
집합으로 가는 함자 $A\longmapsto F(A)$가 적당한 $\mathcal O$에 대하여
$\operatorname{Hom}(\mathcal O,A)$ 꼴로 나타날 수 있기 위한 간단한
필요충분조건을 주기는 쉽다. 여기서는 그 조건들을 자세히 말하지 않겠다.
다만 지금의 경우 이 조건들은 $X_0$에 비자명한 코호몰로지적 조건들을
부과하며, 강연자가 반례를 구성하지는 않았지만 그 조건들이 언제나
성립할 가능성은 낮아 보인다. 한편
$H^0(X_0,\Theta_{X_0/k})=0$이라는 조건은 형식적 모듈라이 스킴의 존재를
위한 충분조건일 듯하다(다만 결코 필요조건은 아니다). 여기서는 특히
단순한 경우의 정리 하나만 진술한다. 해석공간 이론에서 이에 대응하는
결과는 잘 알려져 있으며(cf. Kodaira--Spencer), 이 정리는 앞 절의
결과들을 사용하여 어렵지 않게 증명된다.

\result{정리}{10} $X_0$를 체 $k$ 위에서 고유이고 비특이인 스킴이라 하고,
다음을 가정하자.
\[
  H^0(X_0,\Theta_{X_0/k})=H^2(X_0,\Theta_{X_0/k})=0.
\]
그러면 $X_0$의 형식적 모듈라이 스킴이 존재하며, 이는 정칙 국소환
$\mathcal O$, 즉 $k$ 위의 형식적 멱급수 대수에 대응한다.

이미 지적했듯이, 일반적으로 $\mathcal O$ 위의 형식적 스킴 $\mathfrak X$가
대수화 가능한 것은 아니다. 그러나 $X_0$가 사영적이고
$H^2(X_0,\Osh_{X_0})=0$이면 대수화 가능함이 알려져 있다(정리 4). 예를
들어 $X_0$의 차원이 1이면 그러하다. 이는 주어진 종수의 곡선들을 위한
정수환 위의 모듈라이 스킴을 구성할 수 있으리라는 기대를 어느 정도
뒷받침한다\ldots{}

또한 이 절에서 설명한 것과 같은 방법들은 피카르 다양체의 구성과 연구뿐
아니라 그 밖의 여러 구성에도 적용할 수 있다. 이에 관해서는 곧 다시
돌아오겠다.
\fsection{8}{기본군에의 응용}

앞에서 설명한 기법들을 사용하면 위상 이론을 본보기로 삼아 기본군을
체계적으로 연구할 수 있다. 이 절에서 서술하는 첫 두 정리는 랑--세르의
최근 논문에 있는 결과들을 일반화한 것이다.

$X$를 스킴이라 하자. $X$-스킴 $X'$이 다음 조건들을 만족하면 이를
$X$의 비분기 피복이라 한다.
\begin{enumerate}
  \item[\textit{i.}] $X'$은 $X$ 위에서 유한이다. 즉 $X$ 위의 연접인
    대수층 $\mathcal A=\mathcal A(X')$에 의해 정의된다.
  \item[\textit{ii.}] $\mathcal A$는 $X$ 위에서 국소 자유층이다.
  \item[\textit{iii.}] 모든 $x\in X$에 대하여 몫
    \[
      \mathcal A_x/\mathfrak m_x\mathcal A_x
      =\mathcal A_x\otimes_{\Osh_x}\kappa(x)
    \]
    은 $\kappa(x)$ 위의 분리 대수이다.
\end{enumerate}

\src{210}
이 비분기 피복의 개념(세르와 강연자에 의한 것이다)은 합리적으로 기대할
수 있는 모든 초보적 성질을 가지며, 여기서는 그 목록을 제시하지 않겠다.
고전적 갈루아 이론과 위상 피복의 갈루아 이론을 본보기로 삼은 갈루아
이론이 여기서 생긴다는 것만 말해 두자. 이 이론은 고전적 갈루아 이론을
포함하며, 그 증명들은 후자에 관하여 통상 받아들여지는 증명들보다 오히려
더 간단하다. 정확히 말해, 스킴 $X$의 \emph{기하적 점}이란 대수적으로
닫힌 체 $\Omega$의 스펙트럼 $\xi$에서 $X$로 가는 사상 $a$를 말한다.
이는 $X$의 점 $x=|a|$의 잉여체 $\kappa(x)$의 대수적으로 닫힌 확대체를
주는 것과 같다(이 $x$를 기하적 점 $a$의 \emph{소재점}이라 한다).
$X'$이 $X$의 비분기 피복이면 이에 ``$a$ 위에 있는 $X'$의 기하적
점들''의 집합 $E_a(X')$를 대응시킬 수 있다. 즉 이는 $x$ 위에 있는
$x'\in X'$과 $\kappa(x)$-준동형
$\kappa(x')\longrightarrow\Omega$로 이루어진 순서쌍들의 집합이다.
따라서 $(X,a)$를 고정하면 $X$의 비분기 피복 $X'$들의 범주 $R(X)$에서
유한집합의 범주로 가는 함자 $F(X,a)$를 얻는다. $X$가 연결이면 $R(X)$와
$F(X,a)$로 이루어진 쌍은 적당한 콤팩트 완전 불연결 위상군 $\pi$(즉
유한군들의 사영극한)가 정의하는 유사한 쌍과 동형이 되기 위해 필요한
형식적 성질들을 갖는다. 곧 $\pi$가 연속적으로 작용하는 유한집합 $E$들의
범주 $C(\pi)$와, 이 범주에서 유한집합의 범주로 가는 망각함자
$F(\pi)(E)=E$를 취한다. 더욱이 $(C(\pi),F(\pi))$가 주어진 쌍과
동형이라는 조건은 $\pi$를 정준 동형까지 결정한다. 이 경우 $\pi$를
기하적 점 $a$에서의 연결 스킴 $X$의 \emph{기본군}이라 하며
$\pi_1(X,a)$로 쓴다. $X$가 연결이 아니면 $X$를 $x=|a|$를 포함하는
연결 성분으로 바꾼다. 한편 $X$가 연결이면 $X$의 두 기하적 점
$a',a''$에 대응하는 군 $\pi_1(X,a)$들은 동형이며(그 동형은 내부
자기동형까지 정해진다), 따라서 통상처럼 목적에 가장 알맞게 $a$를
고를 수 있다. 예를 들어 $X$가 기약이라고 가정하고 그 일반점에서
$a$를 고를 수 있다. 물론 $\pi_1(X,a)$는 기점이 지정된 스킴 $(X,a)$에
관한 공변 함자이다. 비분기 피복의 분류에 관한 모든 명제\footnote{인쇄된
본문에는 ``비분리 피복''이라고 되어 있다. 그러나 이 문단은 비분기
피복들의 범주 $R(X)$와 연속적인 $\pi_1(X,a)$-작용을 갖는 유한집합들의
범주 사이의 동치를 다룬다.}는 잘 알려진 사전에 따라 군론의 언어로
옮길 수 있다(다만 여기서는 군들이 위상군이라는 사실을 고려해야 한다).

\src{211}
우리의 목적은 고유 사상 $f:X\longrightarrow Y$에 관하여 섬유공간의
호모토피 완전열에 대응하는 것을 전개하는 데 있다. 고차 호모토피군이
무엇인지 알지 못하므로, 얻는 결과가 불완전할 수밖에 없음은 분명하다.
제3절의 기본정리들을 적용하려면 먼저 체 또는 아르틴 환 위의 스킴에
관한 몇 가지 초보적 보조정리를 명시해야 한다(일반적인 절차에 따라서!).

\result{보조정리}{1} $(X',a')$를, 표시점 $e$를 갖춘 유한집합 $E$
위에서의 $\pi_1(X,a)$의 표현에 대응하는 기점이 지정된 비분기 피복이라
하자. 그러면 정준 준동형
$\pi_1(X',a')\longrightarrow\pi_1(X,a)$은 첫째 군을 $\pi_1(X,a)$ 안의
$e$의 안정자 부분군과 동일시한다(따라서 이 준동형은 단사이다).

\result{보조정리}{2} $X$를 체 $k$ 위의 대수적 스킴, $k'$을 $k$의
순수 비분리 확대체라 하자. 그러면 $X\otimes_k k'$의 모든 비분기 피복은
$X$의 어떤 비분기 피복의 역상(즉 스칼라 확대)으로 얻어지며, 그 피복은
정준 동형까지 결정된다.

특히 이 두 보조정리로부터, $k$의 모든 대수적 확대체 $K$와
$X$의 기하적 점 $a$로 사영되는 $X'=X\otimes_k K$의 모든 기하적 점
$a'$에 대하여 함자적 준동형
$\pi_1(X',a')\longrightarrow\pi_1(X,a)$이 단사임이 따른다.

\result{보조정리}{3} $A$를 아르틴 국소환이라 하고, $X$를 $A$ 위의
완전 스킴으로서 $H^0(X,\Osh_X)=A$인 것이라 하자. $X'$을 $X$의 비분기
피복이라 하고 $A'=H^0(X',\Osh_{X'})$로 놓자. 그러면 $A'$은
$A$-가군으로 유한 생성인 환이다(선험적으로는 $A$ 위에서 분기될 수도
있다). $X_0,X'_0$를 각각 $X,X'$에 딸린 축소 부분스킴이라 하자(이는
$\Osh_X,\Osh_{X'}$ 안의 멱영원들로 이루어진 층으로 나누어 얻는다).
$A/\mathfrak r(A)$의 부분체 $k$로서 $A/\mathfrak r(A)$가 $k$ 위에서
유한인 것을 택하자(따라서 $X_0$는 $k$ 위의 완전 대수적 스킴이고
$X'_0$는 그 비분기 피복이다). 끝으로 $\Omega$를 $k$의 대수적으로 닫힌
확대체라 하고, $X_0\otimes_k\Omega$의 비분기 피복
$X'_0\otimes_k\Omega$를 생각하자. 그러면 다음 두 조건은 동치이다.
\begin{enumerate}
  \item[\textit{i.}] $X'_0\otimes_k\Omega$는
    $X_0\otimes_k\Omega$ 위에서 완전히 분해되어 있다.
  \item[\textit{ii.}] 자연스러운 사상
    $X'\longrightarrow X\otimes_A A'$은 동형이다.
\end{enumerate}
이 조건들이 성립하면 $A'$은 $A$의 비분기 확대이다. 끝으로 $X'$이
연결이면 조건 \textit{(i)}는 겉보기에는 더 약한 다음 조건과 동치이다.
\begin{enumerate}
  \item[\textit{i bis.}] $X'_0\otimes_k\Omega$는
    $X_0\otimes_k\Omega$ 위에서 정칙 단면을 갖는다.
\end{enumerate}
조건 \textit{(ii)}가 성립할 때 $X$의 비분기 피복 $X'$을
\emph{기하적으로 자명하다}고 하겠다.

\src{212}
\result{보조정리}{4} $f:X\longrightarrow Y$를
$f_*(\Osh_X)=\Osh_Y$를 만족하는 고유 사상이라 하자. $a$를 $X$의
기하적 점, $b$를 $Y$ 위의 그 상이라 하자. 그러면
$\pi_1(X,a)\longrightarrow\pi_1(Y,b)$는 전사이다.

실제로 다음을 보이면 충분하다. $Y$의 비분기 피복 $Y'$이 국소 자유인
대수의 층 $\mathcal A$에 대응하고 $X\otimes_Y Y'$이 비연결이면,
$Y'$도 비연결이다. 실제로 $f^*(\mathcal A)$는 서로 영이 아닌 두 환의
층의 직합이고, 따라서 그 직접상도 그러하다. 그런데 이 직접상은 다름
아닌
$\mathcal A\otimes f_*(\Osh_X)=\mathcal A$이다.

\result{보조정리}{5} $X$를 체 $k$ 위의 완전 스킴이라 하자.
$H^0(X,\Osh_X)$가 국소환 $A$이고 $A/\mathfrak r(A)$가 $k$의 순수
비분리 확대체라고 가정하자. $\Omega$를 $k$의 대수적 폐포라 하고
$\overline X=X\otimes_k\Omega$라 하자(이는 연결이다). $\overline X$의
기하적 점 $\overline a$로서 $X$의 기하적 점 $a$로 사상되는 것을 하나
택하면 다음 완전열을 얻는다.
\[
  e\longrightarrow\pi_1(\overline X,\overline a)
  \longrightarrow\pi_1(X,a)
  \longrightarrow\pi_1(k,b)\longrightarrow e
\]
(여기서 $\pi_1(k,b)$는 $\Omega$의 $k$ 위의 갈루아 군이다.)

첫 번째 준동형이 단사라는 사실은 보조정리 1과 2에서 보았고, 가운데에서의
완전성은 보조정리 3에서 따르며, 마지막 준동형의 전사성은 보조정리 4에서
따른다. 이 마지막 사실만이 $A/\mathfrak r(A)$가 $k$의 순수 비분리
확대체라는 가정을 사용한다.

\result{명제}{4} $f:X\longrightarrow Y$를 고유 평탄 사상이라 하고,
모든 $y\in Y$에 대하여
$H^0(f^{-1}(y),\Osh_{f^{-1}(y)})$가 잉여체 $\kappa(y)$ 위의 분리가능
대수라고 하자. 예를 들어 $f^{-1}(y)$가 $\kappa(y)$ 위의 분리가능
스킴, 즉 축약이고 그 기약 성분들에 대응하는 체들이 $\kappa(y)$의
분리가능 확대체이면 이 조건이 성립한다. 그러면 $f_*(\Osh_X)$에
대응하는 $Y$의 피복 $Y'$은 비분기이다.

증명은 정리 2로부터 쉽게 따른다.

\src{213}
이 명제와 보조정리 1을 함께 쓰면, 분리가능 섬유를 갖는 고유 평탄
사상들의 호모토피론적 연구는 사실상 $f_*(\Osh_X)=\Osh_Y$인 경우로
환원된다. 실제로 스타인 인수분해를 사용하여 $Y$를 $Y'$으로 바꾸면 된다.

\result{비고}{} 유한형 평탄 사상으로서 그 섬유들이 분리가능 스킴
(각각 비특이 스킴)인 것을 \emph{분리가능}(각각 \emph{매끄러운})
사상이라 한다. $f$가 평탄하고 $f^{-1}(y)$가 분리가능(각각 비특이)이면,
$f^{-1}(y)$의 어떤 근방 위에서 $f$가 분리가능(각각 매끄러운) 사상이
됨을 보일 수 있다. 같은 결과가 “절대 정규”에 대해서도 성립한다
(“베르티니 정리”).

$f:X\longrightarrow Y$를 다음 조건을 만족하는 고유 사상이라 하자.
\begin{equation}
  \tag{i}
  f_*(\Osh_X)=\Osh_Y.
\end{equation}
또 $X'$을 $X$ 위에서 유한인 스킴이라 하자. $X'\longrightarrow Y$의
스타인 인수분해에 대응하는 $Y$의 피복을 $Y'$이라 하자(정리 5 참조).
$y\in Y$라 하자. 그러면 $y$ 위의 $X'$의 섬유 $F'$의 연결 성분들의
집합은 $y$ 위에 놓인 $Y'$의 점들의 집합과 동일시된다(정리 5). 자연스러운
사상 $X'\longrightarrow X$와
$X'\longrightarrow Y'$로부터 얻어지는 자명한 사상
\begin{equation}
  \tag{*}
  X'\longrightarrow X\times_Y Y'
\end{equation}
을 생각하자.\footnote{인쇄본에서는 두 번째 $X'$의 프라임을 빠뜨렸다.
$Y'$로 가는 사상은 $X'\to Y$의 스타인 인수분해에서 오는 사상이다.}
$Y$의 비분기 피복 $Y''$에 대하여 $X'=X\times_Y Y''$ 꼴이면 이 사상은
언제나 동형이고, 이때 $Y'=Y''$이며 (*)는 항등사상이다. 우리가 구하려는
것은 바로 $X'$이 이와 같은 꼴, 즉 $Y'$이 비분기이고 (*)가 동형이 되기
위한 조건이다. 이를 위하여 $y$에서 $X$의 섬유를 $F$라 하자. 이는
$\kappa(y)$ 위의 고유 스킴이며 $F'$은 $F$의 피복이다($X'$이 $X$ 위에서
비분기이면 이 피복도 비분기이다). $F'_1$을 $y$ 위에 놓인 $Y'$의 한 점
$y'_1$에 대응하는 $F'$의 연결 성분이라 하자. 다음 조건들을 더 가정한다.
\begin{enumerate}
  \item[\textit{(ii)}] $X'\longrightarrow X$는 $F'_1$의 점들에서
    비분기이다. 따라서 $F'_1$은 $F$의 비분기 피복이다.
  \item[\textit{(iii)}] $F'_1$은 $F$의 기하적으로 자명한 피복이다
    (보조정리 3 참조).
\end{enumerate}

\result{정리}{11} 이 조건들 아래에서 $Y'$ 안의 $y'_1$의 열린 근방
$U'$이 존재하여 (*)는 $U'$ 위에서 동형이다. 더 나아가
$Y'\longrightarrow Y$는 $y'_1$에서 비분기이다(그러나 $y$ 위에 놓인
$Y'$의 다른 점 $y'$들에서는 분기될 수 있다).

물론 조건 \textit{(ii)}와 \textit{(iii)}는 정리의 결론이 성립하기
위한 필요조건이기도 하다. 증명은 보조정리 3과 정리 2를 쓰면 쉽게 따른다.

\src{214}
\result{따름정리}{1} 여전히 \textit{(i)}가 성립한다고 가정하자. $X$ 위의
비분기 피복 $X'$이 $Y$의 비분기 피복 $Y'$의 역상과 동형일 필요충분조건은
$X'$이 각 섬유 $f^{-1}(y)$ 위에 기하적으로 자명한 피복을 유도하는 것이다.

정리 11에 따라 이 조건을 만족하는 $Y$의 점들의 집합은 열려 있으므로,
닫힌점 $y$들에 대해서만 확인하면 충분하다\ldots{} 따름정리 1과 동치인
다음 명제도 기록해 두자. (보조정리 4에 따라 전사인) 준동형
$\pi_1(X)\longrightarrow\pi_1(Y)$의 핵은
$\pi_1(\overline{f^{-1}(y)})$들의 $\pi_1(X)$ 안에서의 상들로 생성되는
닫힌 정규부분군이다. 여기서
\[
  \overline{f^{-1}(y)}
  =f^{-1}(y)\otimes_{\kappa(y)}\overline{\kappa(y)},
\]
$\overline{\kappa(y)}$는 $\kappa(y)$의 대수적 폐포를 나타낸다. 모든
섬유에 같은 기점을 택할 수 없으므로,
$\pi_1(\overline{f^{-1}(y)})\longrightarrow\pi_1(X)$인 준동형들은
(일단 $X$의 기점을, 따라서 $Y$의 기점도 택하고 나면) 어쨌든
$\pi_1(X)$의 내부 자기동형사상을 뒤에 합성하는 차이까지만 정해진다는
점에 주의하자.

\result{따름정리}{2} 정리 11의 일반 가정 아래에서, 더 나아가
$Y,X,X'$이 정역 스킴이라고 가정하고 그 유리함수체를 각각 $K,L,L'$이라
하자. 그러면 $L'$ 안에 $K$의 분리 중간확대체 $K'$가 존재하여, $K'$과
$L$은 $K$ 위에서 선형적으로 서로소이고 $L'=LK'$이다(따라서
$L'=L\otimes_K K'$이다).

(보조정리 3의 마지막 부분을 $X$의 기하적 섬유에 적용한다.) 정리 11의
가장 흥미로운 적용 경우는 $f$가 분리 가능한 사상일 때 얻어진다. 이때
$X'$도 $Y$ 위에서 분리 가능하다. 따라서 명제 4에 의해 $Y'$은 $Y$ 위에서
비분기이고, 그러므로 (*)의 우변 $X\times_Y Y'$은 $X$ 위에서 비분기이다.
이에 따라 쉽게 다음을 얻는다.

\result{따름정리}{3} \textit{(i)}에 더하여 $f$가 분리 가능하다고 가정하자.
$X'$을 $X$의 연결인 비분기 피복이라 하자. $X'$\footnote{인쇄본에는
여기서 ``$X$''라고 되어 있으나, 이 판정의 주어는 앞 문장에서 도입한
피복 $X'$이다.}이 $Y$의 비분기 피복 $Y'$의 역상일 필요충분조건은, 한
기하적 섬유 $\overline F=f^{-1}(y)$ 위에 유도된 피복 $\overline F'$이
정칙 단면을 갖는 것이다.

$\overline F'$이 $\overline F$ 위에서 기하적으로 자명하다고 가정할 필요는
없었다는 점에 주의하자(이것은 \textit{a posteriori} 참이지만,
\textit{a priori}에는 훨씬 더 강한 조건이다). 실제로 따름정리 3은 다음
명제와 동치이다.

\src{215}
\result{따름정리}{4} $f:X\longrightarrow Y$를
$f_*(\Osh_X)=\Osh_Y$를 만족하는 고유이고 분리 가능한 사상이라 하자.
$y\in Y$의 기하적 섬유를 $\overline F$라 하고, $\overline F$ 안에서
기하적 점 하나를 택하자. 그러면 사상들
$\overline F\longrightarrow X\longrightarrow Y$에 의해 $X,Y$에서도
기하적 점들이 정해지며, 이들을 $\overline F,X,Y$의 기본군의 기점으로
삼는다. 이 조건들 아래에서 다음 완전열을 얻는다.
\[
  \boxed{
    \pi_1(\overline F)\longrightarrow\pi_1(X)
    \longrightarrow\pi_1(Y)\longrightarrow e
  }.
\]

이로부터 정규성에 관한 어떤 가정도 필요 없는 다음 두 세르--랭 명제를
쉽게 얻는다.

\result{따름정리}{5} $X,Y$를 체 $k$ 위의 연결인 스킴들이라 하자. 축소
스킴 $X_{\mathrm{réd}}$가 $k$ 위에서 분리 가능하다고 가정한다($k$가
완전체이면 이는 자동으로 참이다). 또한 $X$ 또는 $Y$가 $k$ 위에서
고유라고 가정한다.\footnote{1962년 정오표에서 추가된 조건이다.} $X$(각각
$Y$) 안의 기하적 점 $a$(각각 $b$)를 택하면 $X\times_kY$ 안의 기하적 점
$c=(a,b)$와 자연스러운 사상
\[
  \pi_1(X\times_kY,c)\longrightarrow
  \pi_1(X,a)\times\pi_1(Y,b)
\]
을 얻는다(이는 $\pi_1(X\times Y,c)$에서 $\pi_1(X,a)$와
$\pi_1(Y,b)$로 가는 $\pi_1$의 함자적 사상들로부터 유도된다). 이 사상은
단사이고, $k$가 대수적으로 닫혀 있으면 전단사이기까지 하다.

(마지막 경우의 전사성은 거의 자명하다.) 이로부터 세르--랭과 함께 다음을
얻는다.

\result{따름정리}{6} $X$를 대수적으로 닫힌 체 $k$ 위의 연결이고 고유인
대수적 스킴이라 하고, $K$를 $k$의 대수적으로 닫힌 확대체라 하자. 그러면
$X$와 $X\otimes_kK$의 기본군은 서로 같다. 즉 뒤의 스킴의 모든 비분기
피복은 $X$의 한 비분기 피복으로부터 스칼라를 확대하여 얻어지며, 그 피복은
동형을 제외하고 유일하다.

\result{비고}{}

1) 명제 4를 사용하면 따름정리 4의 가정 $f_*(\Osh_X)=\Osh_Y$가 본질적이지
않음을 알 수 있다. 일반적인 경우에는 $\pi_1(Y)$ 뒤에 자명군 $e$를 놓는
대신 다음 열을 이어야 한다.
\[
  \pi_0(\overline F)\longrightarrow\pi_0(X)
  \longrightarrow\pi_0(Y)\longrightarrow e.
\]
이는 대수적 위상수학에서와 같다.

\src{216}
2) 일반적으로는 현재로서
$\pi_1(\overline F)\longrightarrow\pi_1(X)$의 핵에 관해 아무것도 말할 수
없으며, 여기에는 $\pi_2(Y)$가 관여해야 할 것이다. 그러나 $Y$가 국소환
$A$의 스펙트럼이면 아래의 정리 12($A$가 완비일 때 그렇다고 주장하는
정리)를 이용하여 $\pi_1(\overline F)\longrightarrow\pi_1(X)$가 단사임을
증명할 수 있을 것으로 보인다.

정리 11에서는 정리 1과 정리 2만을 사용하였다. 이제 다음의 초등적인
보조정리를 이용하여 정리 3을 적용하겠다.

\result{보조정리}{6} $X$를 스킴, $X_0$를 이에 대응하는 축소 스킴
(즉 멱영원들을 0으로 만든 스킴)이라 하자. 그러면 $X_0$의 모든 비분기
피복 $X'_0$은 $X$의 비분기 피복 $X'$에서 유도되며, $X'$은 표준 동형을
제외하고 결정된다.

순전히 국소적인 성질의 이 보조정리는 여기서 모듈라이 이론의 정리 8과
비슷한 구실을 한다. 이를 존재 정리(정리 3)와 결합하면 다음을 얻는다.

\result{정리}{12} $A$를 잉여체 $k$를 갖는 완비 뇌터 국소환이라 하고,
$X$를 $A$ 위에서 고유인 스킴이라 하자. 그러면
$X_0=X\otimes_A k$의 모든 비분기 피복 $X'_0$은 $X$의 비분기 피복
$X'$에서 유도되며, $X'$은 동형을 제외하고 유일하다.

다시 말하면 다음과 같다.

\result{따름정리}{1} $X_0$의 한 기하적 점을 $X_0$와 $X$의 기본군을 위한
기점으로 택하자. 그러면 표준 준동형
$\pi_1(X_0)\longrightarrow\pi_1(X)$는 동형이다.

이제 $X_0$에 보조정리 5를 적용하자(간단히 하기 위하여
$H^0(X_0,\Osh_{X_0})=k$라고 가정한다). 또한 $A$가 완비이므로 그 비분기
확대들은 잉여체의 비분기 확대들과 대응한다. 즉
$Y=\operatorname{Spec}(A)$라 할 때 $\pi_1(Y)=\pi_1(k)$이다. 이에 따라
다음 완전열을 얻는다.
\[
  e\longrightarrow\pi_1(\overline X_0)
  \longrightarrow\pi_1(X)\longrightarrow\pi_1(Y)\longrightarrow e.
\]

\result{따름정리}{2} $f:X\longrightarrow Y$를 고유 평탄 사상이라 하고,
$y_1$을 $Y$의 한 점, $y_0$를 $y_1$의 한 특수화라 하자. 이에 대응하는
``기하적'' 섬유 $\overline X_1$과 $\overline X_0$를 생각하고,
$\overline X_0$가 분리 가능하고 연결이라고 가정하자(그러면
$\overline X_1$도 같은 조건들을 만족한다). 이때 내부 자기동형을 제외하고
정의되는 군 준동형
$\pi_1(\overline X_1)\longrightarrow\pi_1(\overline X_0)$을 택할 수 있으며,
이 준동형은 전사이다.

\src{217}
이 준동형이 언제나 전단사이기를 기대할 수 있겠지만, 안타깝게도
$\kappa(y_0)$의 표수가 $>0$이면 일반적으로 그렇지 않다. 그러나 아래에서
(적어도 $\overline X_0$가 비특이일 때) 이 준동형의 핵에 대한 상계를 얻을
것이며, 이로부터 $\kappa(y_0)$의 표수가 $0$이면 위 준동형이 전단사임이
따른다(이는 초월적 방법으로도 증명할 수 있는 결과이다). 어쨌든 이미
일반 섬유의 $\pi_1$을 이용하여 특수 섬유의 $\pi_1$에 대한 상계를 얻는다.
예를 들어 표수 $p$의 대수곡선을 표수 $0$의 곡선으로 들어 올릴 수 있다는
사실(정리 9, 따름정리 4)을 이용하면, 초월적 방법으로 다음을 얻는다.

\result{따름정리}{3} $X_0$를 임의의 표수를 갖는 대수적으로 닫힌 체 위의
완비 비특이 곡선의 스킴이라 하고, $g$를 $X_0$의 종수라 하자. 그러면
$\pi_1(X_0)$는 잘 알려진 관계를 만족하는 $2g$개의 위상적 생성원을 갖는다.

잘 알려진 초평면 절단 기법으로부터 다음을 얻는다.

\result{따름정리}{4} $X$를 임의의 표수를 갖는 대수적으로 닫힌 체 위의
비특이 사영 스킴이라 하자. 그러면 $\pi_1(X)$는 유한 개의 위상적 생성원을
갖는다.

준동형 $\pi_1(\overline X_1)\longrightarrow\pi_1(\overline X_0)$의 핵을
구해 보자. 이를 위해 $Y$를 완비 이산 값매김환 $V=\Osh_y$의 스펙트럼이라고
가정해도 된다($y=y_0$). 문제는 다음과 동치이다. $\overline X_1$의 비분기
피복 $\overline X'_1$(원한다면 갈루아 피복이라고 가정해도 된다)가 주어졌을
때, 그것이 어떤 조건 아래 $X_0$의 비분기 피복에서 유도되는지를 결정하는
것이다. \textit{A priori}, 주어진 피복은 $V$의 분수체 $K$의 대수적 폐포
$\overline K$ 안의 유한 부분확대 $K'$에 대하여
$X_1\otimes_K K'$의 비분기 피복 $X'_1$을 스칼라 확대하여 얻어진다.
$\overline X'_1$이 군 $G$를 갖는 갈루아 피복이면, $X'_1$도 군 $G$를 갖는
갈루아 피복으로 택할 수 있다. 그러면 다음이 성립한다.
$X'_1\otimes_{K'}\overline K=\overline X'_1$이 $X_0$의 비분기 피복에서
유도될 필요충분조건은 $\overline K$ 안에 $K'$의 유한 확대 $K''$이 존재하여
$X''_1=X'_1\otimes_{K'}K''$이 $X''\otimes_{V''}K''$ 꼴이 되는 것이다.
여기서 $V''$은 $K''$ 안에서 $V$의 정수적 폐포이고, $X''$은
$X\otimes_VV''$의 비분기 피복이다. 예를 들어 $X_0$가 절대적으로
정규라고 가정하자. 그러면 $X\otimes_VV''$은 정규이다($V''$ 위에서
평탄하고 특수 섬유가 정규이기 때문이다). 그 유리함수체는
\[
  X_1\otimes_KK''
  =(X\otimes_VK)\otimes_KK''
  =X\otimes_VK''
  =(X\otimes_VV'')\otimes_{V''}K''
\]
의 유리함수체 $K''(X_1)$와 같다.

$L''=K''(X'_1)$을 $X'_1\otimes_{K'}K''$의 유리함수체라 하자. 이는
$K''(X_1)$의 유한 분리 확대이고, 위 조건은 $L''$이
$X\otimes_VV''$의 유리함수체의 비분기 확대라는 뜻이기도 하다(즉
$L''$ 안에서 $X\otimes_VV''$을 정규화한 것은 $X\otimes_VV''$ 위에서
비분기이다). 사실 $L''$이 $X\otimes_VV''$의 특수 섬유의 점들에서
비분기임을 확인하는 것으로 충분하다(일반 섬유
$X_1\otimes_KK''$ 위에서는 이미 비분기이기 때문이다). 이제 $X_0$가
비특이라면 Nagata--Zariski의 ``순수성 정리''에 따라, $L''$이
$X\otimes_VV''$의 특수 섬유의 일반점에서의 국소환 $O''$ 위에서
비분기임을 확인하는 것만으로도 충분하다. 이 환은
$K''(X_1)$ 안에서 $X$의 특수 섬유의 일반점에서의 국소환
$O\subset K(X_1)$의 정수적 폐포인 이산 값매김환이다. 따라서 다음을
얻었다.

\src{218}
\result{따름정리}{5} 앞의 조건과 표기 아래에서,
$\overline X_1=X_1\otimes_K\overline K$의 비분기 피복
$X'_1\otimes_{K'}\overline K$가 $\overline X_0$의 비분기 피복에서
유도될 필요충분조건은, $K''(X'_1)$이
$K''(X_1)$의 이산 값매김환 $O''\subset K''(X_1)$ 위에서 비분기가
되게 하는 $\overline K/K'$의 유한 중간 확대체 $K''$가 존재하는 것이다.

이제 $O''$는 $K''(X_1)$ 안에서 이산 값매김환
$O'\subset K'(X_1)$의 정수적 폐포이고($O'$은 $K'(X_1)$ 안에서
$O$의 정수적 폐포이다), $O'$은 $K'$ 안에서 $V$의 정수적 폐포인
$V'$을 포함하며, $V'$의 균등화원 $u$는 $O'$의 균등화원이기도 함을
주목하자. 이제 $X'_1$이 갈루아군 $G$를 갖는 갈루아 피복이고 그 위수
$n$이 $\kappa(y_0)$의 표수 $p$와 서로소라고 하자(이 표수는 $O'$의
잉여체의 표수이기도 하다). 그러면 $K'(X'_1)$은 $O'$에 관하여
``순분기''이고, 이로부터 쉽게 다음 사실이 따른다(``아브얀카르
보조정리''). $O'$의 한 균등화원의 $n$제곱근 $v$를 첨가하면 이 체는
$K'(X_1)(v)$ 안에서 $O'$의 정수적 폐포 위에서 비분기가 된다. 그런데
$v$는 $V'$의 한 균등화원의 $n$제곱근으로 택할 수 있으므로,
따름정리 5의 조건이 성립한다. (아브얀카르 보조정리와 순수성 정리를
사용할 수 있다는 이 가능성은 세르가 내게 팔았다.) 얻은 결과를
표현하기 위하여, 모든 완전 비연결 콤팩트군 $\pi$에 대하여
$\widetilde\pi$를 $\pi$의 $p$-실로우 부분군들로 생성되는 닫힌 부분군에
의한 몫군, 즉 위수가 $p$와 서로소인 $\pi$의 이산 몫군들의 사영극한이라
하자. 이 표기 아래에서 다음을 얻는다.

\result{정리}{13} $f:X\longrightarrow Y$를 고유하고 평탄한 사상이라
하고, $y_1$을 $Y$의 한 점, $y_0$을 $y_1$의 한 특수화라 하자.
$\overline X_0$가 연결이고 비특이라고 가정한다. 그러면 정리 12의
따름정리 2에 나오는 전사 준동형에서 유도되는 준동형
$\widetilde\pi_1(\overline X_1)\longrightarrow
\widetilde\pi_1(\overline X_0)$은 동형이다.

다시 말해 다음과 같다.

\src{219}
\result{따름정리}{1} 갈루아군의 위수가 $\kappa(y_0)$의 표수 $p$와
서로소인 비분기 갈루아 피복들의 분류는 $\overline X_0$와
$\overline X_1$에 대하여 동일하다.

특히 $\kappa(y_0)$의 표수가 0이면, 대수적 방법으로
$\pi_1(\overline X_1)\longrightarrow\pi_1(\overline X_0)$이
전단사임을 얻는다.

끝으로 사용한 기법들은 정리 13보다 더 일반적인 다음 결과도 준다.

\result{정리}{14} $f:X\longrightarrow Y$를 고유하고 매끄러운 사상이라
하고, $D$를 $Y$ 위에서 매끄럽고 모든 점에서 여차원이 1인 $X$의 닫힌
부분스킴이라 하자. $f$의 섬유 $Z=f^{-1}(z)$에 대하여
$Z'=Z-Z\cap D$라 놓고, $\pi_1^t(Z')$를 $\overline{Z\cap D}$ 위에서
순분기인 $Z'$의 비분기 피복들을 분류하는 기본군 $\pi_1(Z')$의 몫군이라
하자. $y_0,y_1$을 정리 13에서와 같이 택하자. 그러면 내부 자기동형까지
정의되는 전사 준동형
\[
  \pi_1^t(\overline X'_1)\longrightarrow\pi_1^t(\overline X'_0)
\]
이 존재하고, 이에 대응하는 준동형
\[
  \widetilde\pi_1^t(\overline X'_1)
  \longrightarrow\widetilde\pi_1^t(\overline X'_0)
\]
은 동형이다.

이로부터 정리 13의 따름정리들과 정리 12의 따름정리 4에 대응하는
변형들이 따른다. 마찬가지로 정리 9의 따름정리 3을 사용하면 초월적
방법으로 다음을 얻는다.

\result{따름정리}{} $X_0$를 임의 표수의 대수적으로 닫힌 체 위의 완비
비특이 곡선의 스킴이라 하고,
$S=(s_i)_{1\le i\le n}$을 $X_0$의 $n$개 원소로 이루어진 유한
부분집합이라 하자. 그러면 $\pi_1^t(X_0-S)$는 $2g+n$개의 위상적 생성원
$x_i,y_i$ ($1\le i\le g$)와 $\sigma_j$ ($1\le j\le n$)를 가지며,
이들은 다음 관계를 만족한다.
\[
  \left(\prod_i x_i y_i x_i^{-1}y_i^{-1}\right)
  \sigma_1\cdots\sigma_n=1,
\]
여기서 $\sigma_j$들은 $s_j$에 대응하는 관성군의 생성원이다. 표수와
서로소인 위수를 갖고 앞의 관계를 만족하는 원소
$x_i,y_i,\sigma_j$들로 생성되는 모든 유한군 $G$에 대하여, 군이
$G$이고 점 $s_j$에서의 관성군이 $\sigma_j$로 생성되는
$X_0-S$의 비분기 갈루아 피복이 존재한다.

$X_0$의 종수가 0이고 $n=3$이면, 적어도 위수가 표수와 서로소인 갈루아
피복에 대해서는 ``세 점 문제''의 해를 얻는다. (여기서는 정리 9가
실제로 필요하지 않으며, 한편 앞의 따름정리를 세 점 문제에서 다루는
특수한 경우로부터 이끌어 낼 수 있을 것으로 보인다.)

\src{220}
\result{비고}{}

1) $X,X_0,X_1$의 일반화된 갈루아 피복(갈루아군이 무한소적일 수도
있는 피복)을 아마도 사용하게 될 더 완전한 연구를 통하여 정리 12의
따름정리 2에서 핵을 복원할 수 있을 것이다. 반면 순분기가 아닌 분기를
허용하는 피복을 연구하는 일은 훨씬 더 어려울 것으로 보인다.

2) 보조정리 6을, 비특이 사영 스킴을 한 초평면 단면을 따라 형식적으로
완비화한 것에 관한 그라우어트의 결과(또는 정리 11 뒤의 비고 2에서
언급했으나 아직 증명되지 않은 정리)와 함께 사용하면, ``추상''
대수기하학에서도 기본군에 관한 고전적인 렙셰츠 정리를 증명할 수 있다.

\begin{center}
  \textbf{참고문헌}
\end{center}

\begin{description}
  \item[{[1]}] Dieudonné (J.) 및 Grothendieck (A.). --- \emph{Éléments de
    géométrie algébrique}, Publications mathématiques de l'Institut des Hautes
    Études Scientifiques (출간 예정).
  \item[{[2]}] Grothendieck (Alexander). --- \emph{The cohomology theory of
    abstract algebraic varieties}. International Congress of Mathematicians
    [1958. Edinburgh] (출간 예정).
  \item[{[3]}] Serre (Jean-Pierre). --- \emph{Faisceaux algébriques
    cohérents}, Annals of Math., Series 2, 제61권, 1955, 197--278쪽.
  \item[{[4]}] Serre (Jean-Pierre). --- \emph{Géométrie algébrique et
    géométrie analytique}, Ann. Institut Fourier Grenoble, 제6권, 1955--1956,
    1--42쪽.
\end{description}

% FGA 정본 인쇄문 보존 장치: 역사적 통합 R26.
\par\smallskip
\begingroup\footnotesize
\noindent\textit{FGA-CANON-ERR-0045. 따름정리 6에서 고유성 가정이 누락된 인쇄본은 C111에 계속 기록되어 있으며, FGA-CANON-DISP-0045가 수정된 현행 진술을 규율한다.}\par
\noindent\textit{FGA-CANON-ERR-0047. $\operatorname{rang}(E)\geq2$ 조건이 누락된 인쇄본은 C113에 계속 기록되어 있으며, FGA-CANON-DISP-0047이 수정된 현행 진술을 규율한다.}\par
\endgroup
